% Generated draft: requires mathematical and source review.
\documentclass[11pt]{article}
\usepackage[margin=25mm]{geometry}
\usepackage{fontspec}
\setmainfont{Latin Modern Roman}
\usepackage{amsmath,amssymb,mathtools,unicode-math}
\setmathfont{Latin Modern Math}
\usepackage{enumitem,needspace,xurl,hyperref}
\setlength{\parindent}{0pt}
\setlength{\parskip}{5pt}
\setlength{\emergencystretch}{3em}
\newcommand{\N}{\mathbb{N}}
\newcommand{\Z}{\mathbb{Z}}
\newcommand{\Q}{\mathbb{Q}}
\newcommand{\R}{\mathbb{R}}
\newcommand{\C}{\mathbb{C}}
\newcommand{\F}{\mathbb{F}}
\newcommand{\A}{\mathbb{A}}
\newcommand{\PP}{\mathbb{P}}
\newcommand{\E}{\mathbb{E}}
\newcommand{\eps}{\varepsilon}
\newcommand{\dd}{\,\mathrm{d}}
\begin{document}
\section*{UnsolvedMath ID: 30000356}
\textbf{Source problem code:} catalogue entry
\subsection{Definitions and mathematical statement}
Write \(\mathbb T=\R/\Z\), let \(v:\mathbb T\to\R\) be the smooth potential, and let \(\omega\) satisfy the source's weak Diophantine condition. In the standard discrete one-frequency convention, the Schr\"odinger operator is
\[
(H_{\lambda,\omega,\theta}u)_n=u_{n+1}+u_{n-1}+\lambda v(\theta+n\omega)u_n
\quad(n\in\Z),
\]
acting on \(\ell^2(\Z)\), with \(\theta\in\mathbb T\) and real coupling \(\lambda\). For \(E\in\R\), define
\[
A_{E,\lambda}(x)=\begin{pmatrix}E-\lambda v(x)&-1\\1&0\end{pmatrix},
\quad A_N(x)=A_{E,\lambda}(x+(N-1)\omega)\cdots A_{E,\lambda}(x),
\]
and the integrated Lyapunov exponent
\[
L(E,\lambda,\omega)=\lim_{N\to\infty}\frac1N\int_{\mathbb T}\log\|A_N(x)\|\,\dd x.
\]
This makes the standard operator and averaged-exponent conventions explicit; it does not assert uniform pointwise phase convergence.

The source conjecture is that for every admissible smooth \(v\) satisfying its transversality condition and every admissible \(\omega\), there exists a finite threshold \(\lambda_0\), which may be enlarged to at least one, such that
\[
\forall\lambda\in\R\ (|\lambda|>\lambda_0)\quad\forall E\in\R,\qquad
L(E,\lambda,\omega)\geq\tfrac14\log|\lambda|>0.
\]
The same threshold works for all energies, including energies not restricted to the spectrum, and both coupling signs are retained. The coefficient \(1/4\) is part of the stated target.

Neither transversality nor weak Diophantine condition is defined in the record; one must not replace them by a Morse condition, a guessed derivative lower bound, or a standard stronger Diophantine inequality. Smooth ordinarily means \(C^\infty\), but any finite differentiability threshold must be recovered from the source. The literal quantifier order allows \(\lambda_0\) to depend on \(v\) and \(\omega\). The assessment calls the target nonperturbative, which might intend frequency-independent threshold control; that stronger dependence requirement is not explicitly supplied and must be checked.

The supplied literature assessment checked 2026-08-21 reports no theorem in the full stated smooth/weak-frequency generality, noting stronger analyticity, regularity, or frequency hypotheses in known results. It does not reproduce those result hypotheses.
\subsection{Short English statement}
For the intended transverse smooth potential and weak Diophantine frequency, ask whether sufficiently large coupling gives an energy-uniform Lyapunov lower bound \(\frac14\log|\lambda|\). The two central hypotheses and any frequency-independent threshold requirement are undefined.
\subsection{Sources}
\begin{itemize}
\item Local UnsolvedMath registry, permanent ID 30000356, catalogue entry, \emph{Lyapunov Positivity for Smooth Quasiperiodic Schr\"odinger Operators}.
\item \emph{Dynamics of Cocycles and One-Dimensional Spectral Theory} (2006), Oberwolfach Reports, pp. 2933--2978, DOI \url{https://doi.org/10.4171/owr/2005/51}.
\item Supplied literature triage checked 2026-08-21 cites only the original report; no separate theorem, transversality formula, or weak Diophantine inequality is reproduced.
\end{itemize}
\end{document}
